% =====================================================================
%  Bhardwaj, Bottini, Fraser-Taliente, Gladden, Gould, Platschorre, Tillim
%  "Lectures on Generalized Symmetries", arXiv:2307.07547v2 (2023) —— 中文译注本
%
%  编译：LuaLaTeX（本机 xdvipdfmx 不稳定，统一用 lualatex，编译 3 遍）
%  用法：powershell -ExecutionPolicy Bypass -File build.ps1
%
%  约定：
%   * 公式编号用 \tag{} 复刻原文编号（如 (3.13)），与原文一一对应（共 745 条，2.1–6.54）；
%   * 图 1–28 为原文矢量图按 400 dpi 截取，图注为译者所译；
%   * % [Tutor Note: ...] 为译者批注，不进入排版，汇总见《译注汇总.md》；
%   * 原文彩色方框分别用 statementbox（蓝）/ definitionbox（绿）/ notebox（粉）排版。
% =====================================================================
\documentclass[11pt,a4paper]{article}

\usepackage[UTF8]{ctex}
\usepackage{amsmath,amssymb,amsthm,mathtools}
\usepackage{bm}
\usepackage{graphicx}
\usepackage[a4paper,left=2.45cm,right=2.45cm,top=2.6cm,bottom=2.7cm]{geometry}
\usepackage[dvipsnames]{xcolor}
\usepackage{enumitem}
\usepackage{booktabs}
\usepackage{array}
\usepackage{multirow}
\usepackage{longtable}
\usepackage{multicol}
\usepackage{fancyhdr}
\usepackage{caption}
\usepackage{tcolorbox}
\tcbuselibrary{breakable,skins}
\usepackage{tikz}
\usepackage[hidelinks,unicode]{hyperref}
\hypersetup{pdftitle={广义对称性讲义（中译注本）},
  pdfauthor={L. Bhardwaj, L. E. Bottini, L. Fraser-Taliente, L. Gladden, D. S. W. Gould, A. Platschorre, H. Tillim; 中译/批注: 本译本},
  pdfsubject={arXiv:2307.07547v2}}

\captionsetup{font=small,labelsep=quad}
\setlength{\emergencystretch}{3em}
\allowdisplaybreaks[1]
\graphicspath{{figures/}{../figures/}}

% ---------- 页眉页脚 ----------
\pagestyle{fancy}
\fancyhf{}
\fancyhead[L]{\small\itshape Lectures on Generalized Symmetries（中译注本）}
\fancyhead[R]{\small\thepage}
\renewcommand{\headrulewidth}{0.4pt}

% ---------- 原文彩色方框 ----------
% 蓝色 Statement 框：\begin{statementbox}{陈述 2.1：标题} … \end{statementbox}
\newtcolorbox{statementbox}[1]{breakable,enhanced,
  colback=blue!3!white, colframe=blue!45!black, colbacktitle=blue!18!white, coltitle=black,
  fonttitle=\bfseries\small, title={#1}, boxrule=0.6pt, arc=2pt,
  left=6pt,right=6pt,top=4pt,bottom=4pt,toptitle=2pt,bottomtitle=2pt, fontupper=\small}

% 绿色 Definition 框：\begin{definitionbox}{定义 3.1：标题} … \end{definitionbox}
\newtcolorbox{definitionbox}[1]{breakable,enhanced,
  colback=green!3!white, colframe=teal!55!black, colbacktitle=green!20!white, coltitle=black,
  fonttitle=\bfseries\small, title={#1}, boxrule=0.6pt, arc=2pt,
  left=6pt,right=6pt,top=4pt,bottom=4pt,toptitle=2pt,bottomtitle=2pt, fontupper=\small}

% 橙色 Example 框：\begin{examplebox}{例 2.1：标题译文} … \end{examplebox}
\newtcolorbox{examplebox}[1]{breakable,enhanced,
  colback=orange!3!white, colframe=orange!55!black, colbacktitle=orange!20!white, coltitle=black,
  fonttitle=\bfseries\small, title={#1}, boxrule=0.6pt, arc=2pt,
  left=6pt,right=6pt,top=4pt,bottom=4pt,toptitle=2pt,bottomtitle=2pt, fontupper=\small}

% 粉色/浅红「带标题」注释框：\begin{notebox}{标题} … \end{notebox}
\newtcolorbox{notebox}[1]{breakable,enhanced,
  colback=red!2!white, colframe=red!40!black, colbacktitle=red!18!white, coltitle=black,
  fonttitle=\bfseries\small, title={#1}, boxrule=0.6pt, arc=2pt,
  left=6pt,right=6pt,top=4pt,bottom=4pt,toptitle=2pt,bottomtitle=2pt, fontupper=\small}

% 无标题注释框（原文个别方框无标题条）：\begin{noteboxplain} … \end{noteboxplain}
\newtcolorbox{noteboxplain}{breakable,enhanced,
  colback=red!2!white, colframe=red!40!black, boxrule=0.6pt, arc=2pt,
  left=6pt,right=6pt,top=4pt,bottom=4pt, fontupper=\small}

% 原文粗体行内小标题 "**Screening.**" → \para{屏蔽。}
\newcommand{\para}[1]{\par\medskip\noindent\textbf{#1}\quad\ignorespaces}

% 手工编号的图题/表题
\newcommand{\figtitle}[2]{\caption*{\small\textbf{图 #1：}#2}}
\newcommand{\tabtitle}[2]{\caption*{\small\textbf{表 #1：}#2}}

% 参考文献条目
\newcommand{\refitem}[2]{\par\hangindent=1.5em\hangafter=1\noindent [#1]~#2}

% ---------- 常用记号（只用这里定义过的宏；不要自造 \newcommand） ----------
\newcommand{\R}{\mathbb{R}}
\newcommand{\C}{\mathbb{C}}
\newcommand{\Z}{\mathbb{Z}}
\newcommand{\N}{\mathbb{N}}
\newcommand{\Q}{\mathbb{Q}}
\newcommand{\cA}{\mathcal{A}}
\newcommand{\cB}{\mathcal{B}}
\newcommand{\cC}{\mathcal{C}}
\newcommand{\cD}{\mathcal{D}}
\newcommand{\cE}{\mathcal{E}}
\newcommand{\cF}{\mathcal{F}}
\newcommand{\cG}{\mathcal{G}}
\newcommand{\cH}{\mathcal{H}}
\newcommand{\cI}{\mathcal{I}}
\newcommand{\cJ}{\mathcal{J}}
\newcommand{\cK}{\mathcal{K}}
\newcommand{\cL}{\mathcal{L}}
\newcommand{\cM}{\mathcal{M}}
\newcommand{\cN}{\mathcal{N}}
\newcommand{\cO}{\mathcal{O}}
\newcommand{\cP}{\mathcal{P}}
\newcommand{\cQ}{\mathcal{Q}}
\newcommand{\cR}{\mathcal{R}}
\newcommand{\cS}{\mathcal{S}}
\newcommand{\cT}{\mathcal{T}}
\newcommand{\cU}{\mathcal{U}}
\newcommand{\cV}{\mathcal{V}}
\newcommand{\cW}{\mathcal{W}}
\newcommand{\cX}{\mathcal{X}}
\newcommand{\cY}{\mathcal{Y}}
\newcommand{\cZ}{\mathcal{Z}}
\newcommand{\fg}{\mathfrak{g}}
\newcommand{\fh}{\mathfrak{h}}
\newcommand{\ft}{\mathfrak{t}}
\newcommand{\fu}{\mathfrak{u}}
\newcommand{\fsu}{\mathfrak{su}}
\newcommand{\fso}{\mathfrak{so}}
\newcommand{\fsp}{\mathfrak{sp}}
\newcommand{\fZ}{\mathfrak{Z}}
\newcommand{\SU}{\mathrm{SU}}
\newcommand{\SO}{\mathrm{SO}}
\newcommand{\Sp}{\mathrm{Sp}}
\newcommand{\U}{\mathrm{U}}
\newcommand{\PSU}{\mathrm{PSU}}
\newcommand{\PSO}{\mathrm{PSO}}
\newcommand{\Spin}{\mathrm{Spin}}
\newcommand{\Pin}{\mathrm{Pin}}
\newcommand{\SymTFT}{\mathrm{SymTFT}}
\newcommand{\dd}{\mathrm{d}}
\newcommand{\ee}{\mathrm{e}}
\newcommand{\ii}{\mathrm{i}}
\newcommand{\Tr}{\operatorname{Tr}}
\newcommand{\tr}{\operatorname{tr}}
\newcommand{\diag}{\operatorname{diag}}
\newcommand{\sgn}{\operatorname{sgn}}
\newcommand{\id}{\mathrm{id}}
\renewcommand{\Re}{\operatorname{Re}}
\renewcommand{\Im}{\operatorname{Im}}
\newcommand{\rk}{\operatorname{rk}}
\newcommand{\rank}{\operatorname{rank}}
\newcommand{\vol}{\operatorname{vol}}
\newcommand{\half}{\tfrac{1}{2}}
\newcommand{\abs}[1]{\left|#1\right|}
\newcommand{\norm}[1]{\left\lVert#1\right\rVert}
\newcommand{\ie}{\textit{i.e.}}
\newcommand{\eg}{\textit{e.g.}}
\newcommand{\cf}{\textit{cf.}}
% 原文的 B 场/规范场常用记号（正体粗体）
\newcommand{\bB}{\mathbf{B}}
\newcommand{\bA}{\mathbf{A}}
\newcommand{\bC}{\mathbf{C}}
